%% Macros de dados do documento
\autor{Angelo Antonio Lima Silveira Filho}
\titulo{RECLin-PT: extração de relações em notas clínicas em português sobre o corpus SemClinBr}
\instituicao{Instituto Federal do Espírito Santo}
\curso{Bacharelado em Sistemas de Informação}
\local{Cachoeiro de itapemirim}
\data{2026}


\tipotrabalho{Trabalho de Graduação}
\preambulo{Trabalho de Conclusão de Curso apresentado à Coordenadoria do Curso de Sistemas de Informação do
Instituto Federal de Educação, Ciência e Tecnologia do Espírito Santo, como requisito parcial para a obtenção do título de Bacharel em Sistemas de Informação.}

\orientador[Orientador][Prof. Dr.]{Cristiano}[da Silveira Colombo]{Instituto Federal do Espírito Santo}
\coorientador[Coorientadora][Profa. Dra.]{Beltrana}[de Tal]{Instituto Federal do Espírito Santo}

\examinadori{Profa. Dra. Fulana de Tal}{Instituto Federal do Espírito Santo}{Examinadora}
\examinadorii{Prof. Dr. Cicrano de Tal}{Instituto Federal do Espírito Santo}{Examinador}

\approvaldate{01}{Agosto}{2026}

% Palavras-chaves
\palavraschave{Extração de relações}[Processamento de linguagem natural clínico][Língua portuguesa][SemClinBr][Aprendizado de máquina]
\keywords{Relation extraction}[Clinical natural language processing][Portuguese language][SemClinBr][Machine learning]

% Ficha catalográfica
\fichabiblioteca{
    (Biblioteca do Campus Cachoeiro de Itapemirim)
}
\fichaautor{Elaborada por XXXXXXXXXXXXXXXXXXXXX – CRB-X/ES - XXX}
% TODO: ficha catalográfica é elaborada pela Biblioteca do campus; os assuntos
% abaixo são uma sugestão a confirmar com o(a) bibliotecário(a).
\fichatags{
    1. Processamento de linguagem natural.
    2. Extração de relações.
    3. Informática em saúde.
    4. Aprendizado de máquina.
    5. Língua portuguesa.
    I. Colombo, Cristiano da Silveira.
    II. Título.
    III. Instituto Federal do Espírito Santo.
}

\cutter{X999y}
\cdd{000.00}
\selectlanguage{brazil}

% Arial (para usar times-new-roman, comente as três linhas abaixo)
\usepackage{helvet}
\renewcommand{\familydefault}{\sfdefault}
\renewcommand{\ABNTEXchapterfont}{\bfseries}
% Comandos para revisar o texto
\usepackage{easyReview}

% ---------------------------------------------------------------------
% AJUSTES DE LAYOUT E AMBIENTE QUADRO (RECLin-PT)
% ---------------------------------------------------------------------
% (1) Reduz estouros de margem (overfull hbox) em parágrafos que contêm
%     termos longos em \texttt (ex.: associated_with, no_relation).
\setlength{\emergencystretch}{3em}

% (2) Permite quebra de linha logo após o sublinhado em nomes monoespaçados,
%     como associated\_with / no\_relation, evitando que estourem a margem.
\makeatletter
\let\reclin@origunderscore\_
\renewcommand{\_}{\reclin@origunderscore\allowbreak}
\makeatother

% (3) Habilita o ambiente flutuante "quadro". O iftex.cls já preparou a
%     "Lista de Quadros" e o contador (\newlistof/\newlistentry), mas deixou
%     o \newfloat comentado para não recriar o contador. Aqui definimos o
%     ambiente reutilizando o contador e a extensão .loq existentes.
\makeatletter
\def\fps@quadro{tbp}
\def\ftype@quadro{16}
\def\ext@quadro{loq}
\def\fnum@quadro{\quadroname\nobreakspace\thequadro}
\newenvironment{quadro}[1][tbp]{\@float{quadro}[#1]}{\end@float}
\makeatother


%% Elementos opcionais
\editardedicatoria{À minha mãe, meu alicerce de força e superação, que esteve ao meu lado mesmo nos dias em que eu já não tinha vontade de continuar. Que jamais me desamparou diante das dificuldades e, com amor, coragem e fé, tornou-se o meu maior exemplo de vida. Tudo o que sou carrega um pouco da sua força.
Aos meus amigos de jornada, que foram âncora e apoio durante todo o processo.}
\editaragradecimentos{% Agradecimentos — elemento opcional (ABNT NBR 14724).
% Referenciado por macros.tex via \editaragradecimentos{...}.
% TODO: redigir antes da defesa.

Agradeço a todos que contribuíram para a realização deste trabalho.
}
\editarepigrafe{% Epígrafe — elemento opcional (ABNT NBR 14724).
% Referenciado por macros.tex via \editarepigrafe{...}.
% Citação seguida de atribuição. TODO: escolher antes da defesa.

\vspace*{\fill}
\begin{flushright}
\textit{``Citação a ser escolhida.''}\\
(Autor, ano)
\end{flushright}
}


%% Listas [Siglas e Simbolos]
\editarlistasiglas{% Lista de siglas — abntex2 environment `siglas` exige pelo menos UM \item,
% caso contrario a compilacao quebra. Lista inicial com siglas que vao
% aparecer no TCC; adicionar/remover conforme o texto evolui.
\begin{siglas}
  \item[ABNT]      Associação Brasileira de Normas Técnicas
  \item[BERT]      \textit{Bidirectional Encoder Representations from Transformers}
  \item[i2b2]      \textit{Informatics for Integrating Biology and the Bedside}
  \item[IC]        Intervalo de Confiança
  \item[IFES]      Instituto Federal do Espírito Santo
  \item[MCC]       \textit{Matthews Correlation Coefficient} (Coeficiente de Correlação de Matthews)
  \item[n2c2]      \textit{National NLP Clinical Challenges}
  \item[NER]       \textit{Named Entity Recognition} (Reconhecimento de Entidades Nomeadas)
  \item[NLP]       \textit{Natural Language Processing} (Processamento de Linguagem Natural)
  \item[PLN]       Processamento de Linguagem Natural
  \item[RE]        \textit{Relation Extraction} (Extração de Relações)
  \item[SemClinBr] \textit{Semantically Annotated Corpus for Portuguese Clinical NLP}
  \item[TF-IDF]    \textit{Term Frequency–Inverse Document Frequency}
  \item[XML]       \textit{Extensible Markup Language}
\end{siglas}
}
\editarlistasimbolos{\input{pre_textuais/simbolos}}